% TEMPLATE-MILC-v3.tex - plantilla maestra UNICA de las guias MILC v3.
%
% La usan las tres familias de guias, cada una con su driver:
%   · Grados 6-11  -> scripts/build-guias-g11.py
%   · Bebras       -> scripts/build-guias-bebras.py
%   · Semillero    -> scripts/build-guias-semillero.py
%
% Antes habia tres copias casi identicas (~400 lineas iguales, 6 distintas):
% cada mejora habia que aplicarla tres veces, y en la practica no se hacia
% (el motor de recursos vivia solo en dos de ellas). Lo que varia entre
% familias esta parametrizado en 6 placeholders que cada driver llena
% (se nombran aqui SIN su sintaxis real: el builder reemplaza por texto plano,
% tambien dentro de los comentarios, y un valor multilinea rompe el preambulo):
%
%   HEADER_IZQ        encabezado izquierdo de pagina
%   HEADER_DER        encabezado derecho de pagina
%   PORTADA_ETIQUETA  rotulo sobre el numero grande (GRADO/MOMENTO/MODULO)
%   PORTADA_SERIE     linea de serie de la portada
%   PORTADA_ID        identificador de la guia en la portada
%   PORTADA_CIRCULO   el circulito de grado (°), o vacio si no aplica
%
% El resto de claves las documenta content/guias/_SCHEMA.md. El script valida
% que no queden placeholders sin reemplazar antes de escribir el archivo final.

\documentclass[11pt,letterpaper]{article}
\usepackage[margin=2.0cm]{geometry}
\usepackage[table]{xcolor}
\usepackage{fontspec}
\usepackage[spanish]{babel}
\usepackage{tikz}
% Librerías TikZ para los diagramas de las guías (bloque `recursos:`):
%  · babel  → babel en español vuelve activos caracteres como «>», y TikZ los usa
%             en sus opciones (>=stealth, ->). Sin esta librería, cualquier
%             diagrama con flechas revienta con «\language@active@arg> has an
%             extra }». Debe ir ANTES de que se dibuje nada.
%  · positioning        → sintaxis «below left=2pt and 3pt».
%  · decorations.pathreplacing → llaves ({brace}) para señalar rangos.
\usetikzlibrary{babel,positioning,decorations.pathreplacing}
\usepackage{graphicx}
% Circuitos: el trabajo de electrónica se hace hoy con micro:bit y sus sensores
% integrados (MakeCode, sin cableado), así que no se necesitan esquemáticos.
% Si algún día entran montajes con protoboard, basta descomentar esta línea:
% los diagramas .tex/.tikz declarados en `recursos:` ya se incrustan solos.
% \usepackage{circuitikz}
\usepackage[most]{tcolorbox}
\usepackage{tabularx}
\usepackage{array}
\usepackage{fancyhdr}
\usepackage{titlesec}
\usepackage[document]{ragged2e}  % bandera a la derecha en todo el cuerpo
\usepackage{enumitem}
\usepackage{microtype}

% Helvetica Neue no trae la flecha U+2192, asi que cada «→» escrito literalmente
% en el YAML se compilaba a NADA, en silencio (462 flechas en 61 guias: rutas del
% tipo «paleta → tipografia → lockup» salian rotas). Mapeamos el caracter a su
% version matematica, que si tiene glifo. Asi el autor puede seguir escribiendo
% «→» comodamente en el YAML.
\usepackage{newunicodechar}
\newunicodechar{→}{\ensuremath{\rightarrow}}
% Mismo problema con los simbolos de marca y vinetas: el «✓» de «una palomita ✓»
% salia como cuadrito vacio en el PDF de todas las guias que lo usaban.
\usepackage{pifont}
\newunicodechar{✓}{\ding{51}}
\newunicodechar{✗}{\ding{55}}
\newunicodechar{✔}{\ding{52}}
\newunicodechar{•}{\textbullet}
\newunicodechar{≥}{\ensuremath{\geq}}
\newunicodechar{≤}{\ensuremath{\leq}}

\setmainfont{Helvetica Neue}
\setsansfont{Helvetica Neue}
\setlength{\parindent}{0pt}
\setlength{\parskip}{6pt}
% Sin particion de palabras: en 6.º-7.º una palabra cortada por guion al final
% de linea («Comuni-cacion») frena la lectura mucho mas que una linea corta.
\hyphenpenalty=10000
\exhyphenpenalty=10000
\pretolerance=2000
\tolerance=3000
\setlength{\emergencystretch}{3em}
\linespread{1.10}
\renewcommand{\arraystretch}{1.25}
\setlist{leftmargin=5mm,itemsep=1mm,topsep=1mm}
\newcolumntype{Y}{>{\RaggedRight\arraybackslash}X}

\definecolor{milcMagenta}{HTML}{E6007E}
\definecolor{milcVino}{HTML}{5A0038}
\definecolor{milcTurquesa}{HTML}{0093A5}
\definecolor{milcVerde}{HTML}{58A923}
\definecolor{milcMostaza}{HTML}{E5B400}
\definecolor{milcOcre}{HTML}{B86600}
\definecolor{milcNegro}{HTML}{241816}
\definecolor{milcGris}{HTML}{F7F7F4}
\definecolor{milcLinea}{HTML}{E8E2DD}
\definecolor{milcGradePrimary}{HTML}{0D9488}
\definecolor{milcGradeDark}{HTML}{115E59}
\definecolor{milcGradeSoft}{HTML}{CCFBF1}
\definecolor{milcGradeAccent}{HTML}{CFE7FF}
\definecolor{milcGradeText}{HTML}{FFFFFF}
\color{milcNegro}

\pagestyle{fancy}
\fancyhf{}
\lhead{\small Territorio interior · Educación socioemocional}
\rhead{\small Grado 6° · Momento 8 · MILC v3}
\cfoot{\small\thepage}
\renewcommand{\headrulewidth}{0pt}

\newtcolorbox{titlebox}[1]{
  enhanced,colback=#1,colframe=#1,arc=7mm,boxrule=0pt,
  left=7mm,right=7mm,top=3mm,bottom=3mm,
  before skip=8pt,after skip=8pt,
  fontupper=\sffamily\bfseries\Large\color{white}
}

\newtcolorbox{softbox}[3]{
  enhanced,colback=#2,colframe=#1,borderline west={4pt}{0pt}{#1},
  arc=2mm,boxrule=.4pt,left=4mm,right=4mm,top=3mm,bottom=3mm,
  before skip=6pt,after skip=8pt,title={#3},coltitle=white,
  colbacktitle=#1,fonttitle=\sffamily\bfseries,fontupper=\RaggedRight,
  attach boxed title to top left={xshift=3mm,yshift=-2mm},
  boxed title style={arc=2mm, boxrule=0pt}
}

\newtcolorbox{drawbox}[2]{
  enhanced,colback=white,colframe=#1,arc=4mm,boxrule=1pt,
  left=5mm,right=5mm,top=4mm,bottom=4mm,before skip=8pt,after skip=8pt,
  title={#2},coltitle=white,colbacktitle=#1,fonttitle=\sffamily\bfseries,
  fontupper=\RaggedRight,
  attach boxed title to top left={xshift=4mm,yshift=-2mm},
  boxed title style={arc=3mm, boxrule=0pt}
}

\newtcolorbox{infoband}[1]{
  enhanced,colback=#1!8,colframe=#1!50,arc=2mm,boxrule=.5pt,
  left=4mm,right=4mm,top=2mm,bottom=2mm,before skip=4pt,after skip=4pt,
  fontupper=\small\RaggedRight
}

% Puente narrativo entre secciones (contrato editorial Conéctate).
% Da continuidad: "Acabas de X. Ahora vas a Y porque Z."
% Visualmente: banda izquierda en color del grado, texto en itálica pequeña.
\newtcolorbox{puentebox}{
  enhanced,colback=milcGradeSoft,colframe=milcGradeDark,
  borderline west={3pt}{0pt}{milcGradeDark},
  arc=0mm,boxrule=0pt,left=5mm,right=4mm,top=2.5mm,bottom=2.5mm,
  before skip=4pt,after skip=8pt,
  fontupper=\itshape\small\color{milcNegro}\RaggedRight
}
% La flecha va en modo matematico a proposito: Helvetica Neue NO tiene el glifo
% U+2192, asi que \textrightarrow se compilaba a NADA («Missing character: There
% is no → in font Helvetica Neue») y el puente salia sin flecha. En modo
% matematico la flecha viene de la fuente matematica y si se dibuja.
\newcommand{\puente}[1]{%
  \begin{puentebox}%
    \textcolor{milcGradeDark}{$\rightarrow$}\hspace{2mm}#1%
  \end{puentebox}%
}

\newcommand{\linead}[1]{\noindent\color{milcLinea}\rule{#1}{0.5pt}\color{milcNegro}}
\newcommand{\respuesta}[1]{\par\vspace{2mm}\foreach \n in {1,...,#1}{\noindent\color{milcLinea}\rule{\linewidth}{0.5pt}\par\vspace{5mm}}\color{milcNegro}}
\newcommand{\checkbox}{\textcolor{milcGradeDark}{\fbox{\phantom{x}}}\hspace{5pt}}

% ─── Listas aireadas para las guías socioemocionales ────────────────────
% Componer las verificaciones como «\checkbox texto\\» deja el texto que salta
% de línea debajo del cuadrito y apelmaza el bloque. Estos entornos alinean la
% continuación y airean los ítems. Los usa Territorio Interior; cualquier
% familia puede hacerlo.
\newlist{checklista}{itemize}{1}
\setlist[checklista]{
  label=\textcolor{milcGradeDark}{\fbox{\phantom{x}}},
  leftmargin=9mm, labelsep=3.2mm, itemsep=2.4mm,
  topsep=2.5mm, parsep=0pt, partopsep=0pt, before=\RaggedRight
}
\newlist{pasoslista}{enumerate}{1}
\setlist[pasoslista]{
  label=\textcolor{milcGradeDark}{\sffamily\bfseries\arabic*.},
  leftmargin=8mm, labelsep=3mm, itemsep=2.8mm,
  topsep=1mm, parsep=0pt, partopsep=0pt, before=\RaggedRight
}

\newcommand{\phasecard}[4]{
\begin{tcolorbox}[enhanced,colback=#3,colframe=#2,arc=3mm,boxrule=.5pt,height=3.05cm,valign=center,halign=center,left=1.5mm,right=1.5mm,top=2mm,bottom=2mm]
{\sffamily\bfseries\fontsize{22}{24}\selectfont\textcolor{#2}{#1}}\par
{\sffamily\bfseries\small #4}
\end{tcolorbox}
}

% ── Piezas del rediseno 2026 (legibilidad 6.º-11.º) ───────────────────
% Banda de estacion: reemplaza el titlebox macizo. Titulo a la izquierda,
% dato practico (tiempo/modalidad) a la derecha, en una sola linea.
\newcommand{\milcestacion}[2]{%
\begin{tcolorbox}[enhanced,colback=#1,colframe=#1,arc=2mm,boxrule=0pt,
  left=5mm,right=5mm,top=2.6mm,bottom=2.6mm,before skip=11pt,after skip=7pt]
{\sffamily\bfseries\fontsize{15}{18}\selectfont\color{white}#2}
\end{tcolorbox}}

% Tarjeta de paso numerada: sustituye las tablas «Pilar / Aplicacion», que
% eran formato de consulta docente, no de estudiante. El numero va en un
% circulo sobre el borde para que la secuencia se lea de un vistazo.
\newcommand{\milcpaso}[3]{%
\begin{tcolorbox}[enhanced,colback=white,colframe=milcLinea,arc=1.5mm,boxrule=.9pt,
  left=14mm,right=4mm,top=3mm,bottom=3mm,before skip=4pt,after skip=4pt,
  fontupper=\RaggedRight,
  overlay={\node[anchor=north west,circle,fill=#2,text=white,inner sep=0pt,
    minimum size=7.4mm,font=\sffamily\bfseries\small]
    at ([xshift=3.6mm,yshift=-3.2mm]frame.north west) {#1};}]
#3
\end{tcolorbox}}

% Casilla marcable de tamano util con lapiz (la anterior era \fbox{\phantom{x}},
% de alto variable segun el texto que la seguia).
\newcommand{\milcmarca}{\framebox[3.8mm]{\rule{0pt}{3.4mm}}}

% Fila marcable de lista suelta (checks de la estacion 1).
\newcommand{\milccheckrow}[1]{%
\par\vspace{1.8mm}\noindent
\makebox[6.4mm][l]{\raisebox{-.55mm}{\textcolor{milcNegro}{\milcmarca}}}%
\parbox[t]{\dimexpr\linewidth-6.4mm\relax}{\RaggedRight #1}\par}

% Celda marcable dentro de la rubrica (tabla de autoevaluacion). Misma
% sangria francesa, pero pensada para vivir dentro de una columna X.
\newcommand{\milcopcion}[1]{%
\makebox[5.2mm][l]{\raisebox{-.55mm}{\textcolor{milcNegro}{\milcmarca}}}%
\parbox[t]{\dimexpr\linewidth-5.2mm\relax}{\RaggedRight #1}}

% ── Recursos visuales de la guía ──────────────────────────────────────
% Declarados en el bloque `recursos:` del YAML y emitidos por el builder.
% \guiaFigura   → imagen rasterizada o PDF (foto, carta celeste, montaje).
% \guiaDiagrama → fuente TikZ/circuitikz (.tex) incrustada como vector:
%                 es la vía para circuitos y esquemáticos editables.
% #1 = ruta relativa al .tex · #2 = pie (puede ir vacío)
\newcommand{\guiaFigura}[2]{%
\begin{center}
\includegraphics[width=0.88\linewidth,height=0.38\textheight,keepaspectratio]{#1}\\[1.5mm]
{\footnotesize\itshape\textcolor{milcNegro!75}{#2}}
\end{center}
}

\newcommand{\guiaDiagrama}[2]{%
\begin{center}
\input{#1}\\[1.5mm]
{\footnotesize\itshape\textcolor{milcNegro!75}{#2}}
\end{center}
}

\begin{document}
\thispagestyle{empty}

% ── Portada ───────────────────────────────────────────────────────────
% Antes: un rectangulo de color a pagina completa. Se veia bien en pantalla,
% pero en la fotocopiadora del colegio se comia el toner y salia gris sucio.
% Ahora la banda ocupa el tercio superior y el resto va en blanco.
\begin{tikzpicture}[remember picture,overlay]
  \path[fill=milcGradePrimary]
    (current page.north west) rectangle ([yshift=-10.6cm]current page.north east);

  \node[anchor=north west,text=milcGradeText] at ([xshift=2.0cm,yshift=-1.55cm]current page.north west)
    {\sffamily\bfseries\fontsize{12}{14}\selectfont MOMENTO};
  \node[anchor=north west,text=milcGradeText,opacity=.85] at ([xshift=2.0cm,yshift=-2.35cm]current page.north west)
    {\sffamily\bfseries\fontsize{8.5}{10}\selectfont TERRITORIO INTERIOR · SOCIOEMOCIONAL};
  \node[anchor=north east,text=milcGradeText,opacity=.85,align=right] at ([xshift=-2.0cm,yshift=-1.55cm]current page.north east)
    {\sffamily\bfseries\fontsize{9}{12}\selectfont I.E. Sor María Juliana\\Cartago, Valle del Cauca};

  \node[anchor=west,text=milcGradeText] (gradonum) at ([xshift=1.85cm,yshift=-7.1cm]current page.north west)
    {\sffamily\bfseries\fontsize{112}{112}\selectfont 8};
  % El circulito de grado (°) solo aplica a las guias de grado («11°»); en
  % Bebras y Semillero el numero grande es un momento/modulo, no un grado.
  % Se posiciona relativo al nodo `gradonum`, no en coordenadas absolutas.
  

  \node[anchor=west,text=milcGradeText,text width=11.4cm,align=left] at ([xshift=1.5cm,yshift=0.5cm]gradonum.east)
    {
     {\sffamily\bfseries\fontsize{9}{11}\selectfont\MakeUppercase{Territorio interior · Socioemocional}}\\[3mm]
     {\sffamily\bfseries\fontsize{21}{25}\selectfont\setlength{\baselineskip}{27pt}La atención secuestrada\\[4pt]quién decide\\[4pt]dónde miras}\\[3.5mm]
     {\sffamily\bfseries\fontsize{15}{18}\selectfont Grado 6° · Momento 8}
    };
\end{tikzpicture}

\vspace*{9.4cm}
\vfill

{\sffamily\bfseries\fontsize{8}{10}\selectfont\textcolor{milcGradeDark}{LO QUE TE LLEVAS HOY}}

\begin{tcolorbox}[enhanced,colback=white,colframe=milcNegro,arc=2mm,boxrule=1.6pt,
  left=5mm,right=5mm,top=4mm,bottom=4mm,before skip=3pt,after skip=12pt,
  fontupper=\RaggedRight\fontsize{11}{15.5}\selectfont]
Tu mapa de atención: dónde se te fue el tiempo ayer, cuáles enganches reconociste en las aplicaciones que usas, y tres decisiones concretas —con hora y lugar— sobre cuándo la pantalla no entra.
\end{tcolorbox}

\begin{tcolorbox}[enhanced,colback=milcGris,colframe=milcGris,arc=2mm,boxrule=0pt,
  left=5mm,right=5mm,top=3.5mm,bottom=3.5mm,after skip=12pt]
{\sffamily\bfseries\fontsize{8}{10}\selectfont\textcolor{milcNegro!55}{ESTA GUÍA ES DE}}\\[3mm]
\begin{tabularx}{\linewidth}{X p{3.2cm} p{3.2cm}}
\linead{\linewidth} & \linead{3.0cm} & \linead{3.0cm}\\[-1mm]
{\fontsize{8.5}{10}\selectfont\textcolor{milcNegro!55}{Nombre completo}} &
{\fontsize{8.5}{10}\selectfont\textcolor{milcNegro!55}{Grupo}} &
{\fontsize{8.5}{10}\selectfont\textcolor{milcNegro!55}{Fecha}}\\
\end{tabularx}
\end{tcolorbox}

\vfill

\noindent\textcolor{milcLinea}{\rule{\linewidth}{1pt}}\\[2mm]
{\sffamily\bfseries\fontsize{9.5}{12}\selectfont Autor: Dr. Álvaro Cárdenas Orozco}\\[1mm]
{\sffamily\fontsize{8.5}{11}\selectfont\textcolor{milcNegro!55}{Metodología MILC · Apertura ancestral · Atención y pantallas · Triángulo Dussel-Estoicismo-Floridi}}

\newpage

\section*{Apertura: La atención secuestrada: quién decide dónde miras}

\begin{softbox}{milcGradeDark}{milcGradeSoft}{Saber ancestral}
Entre los pueblos del Cauca ---nasa, kokonuko--- hay un lugar que enseña sin ser un salón: \textbf{la tulpa}. La palabra viene del quichua y significa \emph{piedra del fogón}: son \textbf{tres piedras} sobre las que se pone la olla, y en nasa yuwe se llama \textbf{ipx kweht}. Las tres piedras no se escogen de cualquier manera ---los mayores las traen con análisis espiritual--- y representan a la familia: padre, madre e hijos. \textbf{Alrededor de ese fuego la familia se reúne, y ahí se enseña la cultura y la lengua}: no con cuaderno, sino \emph{mirando y escuchando}, durante horas. Es un espacio educativo tan real que en los resguardos los cabildos y las guardias estudiantiles se posesionan alrededor de la tulpa. \textbf{Fíjate en cómo funciona:} hay \emph{un} fuego, y todos miran hacia él. La atención está en un solo lugar, es compartida, y dura. \emph{Nadie tiene que pelearla, porque nadie más la está pidiendo.} \textbf{Ahora mira una sala de tu casa a las ocho de la noche:} cuatro personas, cuatro luces distintas, cuatro atenciones tiradas en cuatro direcciones. No es que la gente sea peor que antes. \textbf{Es que ahora hay muchos fuegos, y algunos fueron diseñados para llamarte.}
\end{softbox}

\begin{softbox}{milcTurquesa}{milcGris}{Saber contemporáneo}
En \textbf{1971} ---antes de internet, antes del celular--- el economista \textbf{Herbert Simon} escribió una frase que explica el mundo de hoy: \emph{«lo que la información consume es evidente: consume la atención de quienes la reciben. De ahí que una riqueza de información cree una pobreza de atención»} (Simon, 1971). Es decir: cuando hay infinito para ver, \textbf{lo escaso ya no es la información: eres tú}. Y donde algo es escaso, aparece alguien dispuesto a comprarlo. Por eso muchas aplicaciones son gratis: \textbf{no pagas con plata, pagas con minutos}, y esos minutos se venden. Para conseguirlos, están \textbf{diseñadas} con mecanismos que no son casuales ---el video que arranca solo, la lista que no se acaba nunca, el rojo del aviso, la racha que se pierde si no entras hoy, el «lo vio a las 8:03»---. \textbf{Nada de eso es tu falta de voluntad}: es un trabajo de ingeniería hecho por gente muy bien pagada, y funciona en adultos también. \emph{Por eso esto no se resuelve con más fuerza de voluntad, sino con decisiones tomadas antes ---y con saber ver el enganche cuando aparece---.}
\end{softbox}

\begin{softbox}{milcMostaza}{milcGris}{Pregunta-puente}
\textit{Alrededor de la tulpa hay un solo fuego y todos miran hacia él \textbf{porque así lo decidieron}. \textbf{¿Cuáles de los fuegos que miraste ayer los elegiste tú}\ldots y cuáles te eligieron a ti?}
\end{softbox}

\begin{softbox}{milcVerde}{milcGris}{Saber y saber hacer}
Al terminar podrás: (1) \textbf{identificar} en qué se te fue el tiempo ayer, con datos y no con impresiones; (2) \textbf{explicar} por qué la atención es limitada, qué es un enganche y por qué saltar de una cosa a otra cuesta; (3) \textbf{analizar} tres aplicaciones que uses y nombrar los mecanismos concretos con que buscan retenerte; (4) \textbf{crear} tres reglas propias ---con hora y lugar--- sobre cuándo la pantalla no entra, y elegir tu propio fuego.
\end{softbox}



\small
\begin{tabularx}{\textwidth}{>{\bfseries}p{3.0cm}Y}
\rowcolor{milcGradePrimary!12}Autor & Dr. Álvaro Cárdenas Orozco\\
DBA & Comprende los fenómenos que afectan a su territorio, regula sus emociones ante ellos y acompaña a otros con empatía (transversal 6°--11°, articulado con la política «Escuela, territorio de vida» del MEN, Resolución 006519 de 2025, y la Ley 1620 de 2013)\\
Referentes & Primera ayuda psicológica (OMS, 2011/2012) · Directrices IASC (2007) · USGS y Servicio Geológico Colombiano · Pueblo nasa y la minga (CRIC) · Dussel · Estoicismo · Floridi\\
Producto final & Tu mapa de atención: dónde se te fue el tiempo ayer, cuáles enganches reconociste en las aplicaciones que usas, y tres decisiones concretas —con hora y lugar— sobre cuándo la pantalla no entra.\\
\end{tabularx}
\normalsize

\puente{Ya sabes nombrar lo que sientes, frenar antes de reaccionar y usar tu aire. Hoy te falta lo que pelea contra todo eso: \textbf{lo que se lleva tu atención sin pedirte permiso}.}

\milcestacion{milcGradeDark}{Ruta MILC: así vamos a aprender}
\begin{tabularx}{\textwidth}{>{\centering\arraybackslash}X>{\centering\arraybackslash}X>{\centering\arraybackslash}X>{\centering\arraybackslash}X}
\phasecard{1}{milcVerde}{milcGris}{Escucha\\[-1mm]\small Observo, escucho y pregunto.} &
\phasecard{2}{milcTurquesa}{milcGris}{Sistematización\\[-1mm]\small Reconozco y decido.} &
\phasecard{3}{milcMagenta}{milcGris}{Praxis\\[-1mm]\small Construyo y aplico.} &
\phasecard{4}{milcVino}{milcGris}{Evaluación\\[-1mm]\small Reflexión liberadora.}
\end{tabularx}


\puente{Primero los datos. No lo que crees que usas: lo que el propio teléfono registra. Duele un poco y es el punto de partida.}

\milcestacion{milcVerde}{Estación 1 de 3 · Escucha}

\begin{softbox}{milcVerde}{milcGris}{Escena para escuchar}
\textbf{Actividad 1 · IDENTIFICA --- El ayer, con datos.} \textbf{Esta página es tuya.} Si tienes celular, abre el registro de \emph{tiempo de uso} ---todos los teléfonos lo traen--- y anota tres cosas de \textbf{ayer}: cuánto tiempo en total, las tres aplicaciones que más tiempo se llevaron, y cuántas veces desbloqueaste el teléfono. Si no tienes celular o no lo traes, hazlo con lo que sí: televisión, computador, consola, el celular de la casa; y si nada aplica, hazlo de memoria \textbf{por horas del día}. \textbf{Luego, la pregunta incómoda:} de ese tiempo, ¿cuánto \textbf{elegiste} ---te sentaste a ver algo que querías ver--- y cuánto \textbf{simplemente pasó} ---abriste «un momentico» y levantaste la cabeza cuarenta minutos después---? Escribe los dos números, aunque sean aproximados. \emph{Nadie va a compararlos con los de nadie: el punto no es quién usa más, es cuánto de lo tuyo lo decidiste tú.}
\end{softbox}

{\sffamily\bfseries\fontsize{8}{10}\selectfont\textcolor{milcNegro!55}{MARCA CUANDO LO HAYAS HECHO}}
\milccheckrow{Anotaste el tiempo total de ayer, las tres aplicaciones que más se llevaron y cuántas veces desbloqueaste.}
\milccheckrow{Separaste el tiempo que elegiste del tiempo que simplemente pasó, con dos números.}
\milccheckrow{Identificaste en qué momento del día se te fue más tiempo sin darte cuenta.}

\begin{infoband}{milcVerde}


\textbf{En tu cuaderno (Actividad 1 · IDENTIFICA).} Título: «Actividad 1. Mi ayer, con datos». \textbf{Formato:} los tres datos + los dos números (elegido / simplemente pasó) + la franja del día donde más se te fue. \textbf{Extensión:} ocho renglones. \textbf{Sabes que terminaste cuando} tienes \textbf{números} y no frases como «uso harto». Esta página es tuya: \textbf{nadie más tiene que leerla si tú no quieres}.
\end{infoband}

\puente{Ya tienes tus números. Ahora entiende por qué pasa, para que dejes de creer que es solo culpa tuya.}

\milcestacion{milcTurquesa}{Fase 2 · Sistematización: entender el enganche}

\begin{softbox}{milcTurquesa}{milcGris}{Cuatro claves sobre la atención}
Cuatro claves para entender por qué esto pasa, por qué le pasa también a los adultos y por qué la solución no es «tener más voluntad».
\end{softbox}

\milcpaso{1}{milcTurquesa}{\textbf{Tu atención es limitada, y por eso vale.} No puedes atender a todo: solo a una cosa a la vez, con algo de ruido alrededor. Simon lo dijo en 1971, cuando no existía nada de lo que hoy tienes en el bolsillo: \textbf{«una riqueza de información crea una pobreza de atención»}. Cuando lo que sobra es contenido, \textbf{lo escaso eres tú}: tus minutos, tus ojos. \emph{Y todo lo escaso tiene precio.} Por eso muchas aplicaciones son gratis: no pagas con plata, pagas con tiempo, y ese tiempo se vende a quien pone los anuncios.}
\milcpaso{2}{milcTurquesa}{\textbf{Los enganches están diseñados, y tienen nombre.} No son casualidades ni «así es la app». Reconócelos: el \textbf{reproducir automático} (el siguiente video empieza antes de que decidas), el \textbf{desplazamiento infinito} (la lista nunca termina, así que nunca hay un final natural para parar), la \textbf{notificación en rojo} (el color que el ojo busca primero), la \textbf{racha} (te castiga por no entrar hoy), el \textbf{«visto a las 8:03»} (te obliga a responder ya), la \textbf{recompensa impredecible} (a veces hay algo buenísimo, a veces nada ---y esa incertidumbre engancha más que el premio seguro---). \emph{Nombrar el mecanismo le quita la mitad de la fuerza, igual que nombrar una emoción.}}
\milcpaso{3}{milcTurquesa}{\textbf{Saltar cuesta, aunque no lo sientas.} Cuando estás estudiando y miras el celular «un segundo», no vuelves de inmediato a donde estabas: la cabeza tiene que reconstruir dónde iba, y eso toma más de lo que parece. \textbf{Por eso una hora con quince interrupciones no es una hora de estudio}, es un montón de arranques. \emph{Y explica algo que te habrá pasado: sentarte dos horas y sentir que no rindió nada. No estabas flojo: estabas saltando.}}
\milcpaso{4}{milcTurquesa}{\textbf{Un solo fuego.} La tulpa funciona porque hay \textbf{uno}, y todos miran hacia allá; por eso la conversación dura y lo que se enseña se pega. \textbf{La versión práctica para ti:} decide cuál es tu fuego en cada momento ---la tarea, la conversación, el partido, la comida--- y \textbf{apaga los otros mientras dure}. No para siempre: mientras dure. \emph{Poner el teléfono boca abajo y en otro cuarto no es un castigo: es escoger a qué fuego estás mirando.}}

\begin{drawbox}{milcTurquesa}{Los seis enganches, para reconocerlos en el acto}
\begin{pasoslista}
\item \textbf{Reproducción automática:} el siguiente empieza solo. \emph{Antídoto: apágala en ajustes.}
\item \textbf{Desplazamiento infinito:} nunca hay final. \emph{Antídoto: ponte un final tú ---un temporizador---.}
\item \textbf{Notificación en rojo:} el ojo va primero al rojo. \emph{Antídoto: apaga las notificaciones que no sean de personas.}
\item \textbf{Racha:} te castiga por no entrar. \emph{Antídoto: recuerda que la racha es de la app, no tuya.}
\item \textbf{«Visto a las 8:03»:} presiona a responder ya. \emph{Antídoto: responder no es una obligación inmediata.}
\item \textbf{Recompensa impredecible:} a veces hay algo buenísimo. \emph{Antídoto: es el mecanismo de las máquinas de apuestas; saberlo ya ayuda.}
\end{pasoslista}
\end{drawbox}

\begin{softbox}{milcMostaza}{milcGris}{Errores comunes y su corrección}
\textbf{«Es falta de voluntad mía»:} está diseñado por equipos grandes para funcionarle a cualquiera, incluidos los adultos. \textbf{«Yo controlo, solo miro un momentico»:} mide el momentico con el reloj una vez y compara. \textbf{Confundir descansar con desconectarse:} media hora de videos rápidos no descansa la cabeza igual que caminar o dormir. \textbf{Prohibir todo:} las reglas imposibles se rompen en tres días y dejan la sensación de haber fallado. \textbf{Creer que solo pasa con el celular:} la televisión encendida de fondo también parte la atención. \textbf{Juzgar a los demás:} en tu casa también hay adultos enganchados; esto no es para señalar, es para entender.
\end{softbox}

\begin{infoband}{milcTurquesa}


\textbf{En tu cuaderno (Actividad 2 · EXPLICA).} Título: «Actividad 2. Cómo me enganchan». \textbf{Formato:} las tres aplicaciones que más tiempo te llevaron ayer y, al lado de cada una, \textbf{cuál de los seis enganches} usa contigo y cómo lo notas. \textbf{Extensión:} media página. \textbf{Sabes que terminaste cuando} puedes explicar por qué esto no es «falta de voluntad» sin quedar como excusa.
\end{infoband}

\puente{Ya reconoces el enganche. Es hora de decidir en frío, porque en caliente ---con el video corriendo--- no se decide nada.}

\milcestacion{milcMagenta}{Fase 3 · Praxis: recuperar la mirada}

\begin{softbox}{milcMagenta}{milcGris}{Cómo se arma el mapa de atención}
Las decisiones sobre pantallas hay que tomarlas \textbf{en frío}, igual que la respiración: con el video corriendo no se decide nada. Tres reglas concretas, y un fuego elegido.
\end{softbox}

\milcpaso{1}{milcMagenta}{\textbf{Medir el ayer (ya lo hiciste).} Tus números de la Actividad 1 son el punto de partida. \emph{Sin ellos, dentro de un mes no vas a saber si algo cambió ---y creer que cambió no es lo mismo que saberlo---.}}
\milcpaso{2}{milcMagenta}{\textbf{Cazar enganches (10 min, en parejas).} Abran juntos una aplicación que ambos usen y \textbf{búsquenle los seis enganches}, uno por uno, señalándolos en la pantalla: dónde está el rojo, dónde arranca solo, dónde no termina nunca. \emph{Verlos con alguien al lado cambia la experiencia: lo que parecía «la app» se vuelve un diseño con intenciones.}}
\milcpaso{3}{milcMagenta}{\textbf{Elegir tres reglas (12 min).} Concretas, con \textbf{hora y lugar}, y posibles en tu casa de verdad. Ejemplos: \emph{«el celular no entra a la mesa»}; \emph{«mientras hago tareas, queda en otro cuarto, boca abajo»}; \emph{«apago la reproducción automática hoy mismo»}; \emph{«después de las 9 de la noche queda cargando fuera del cuarto»}; \emph{«las notificaciones que no son de personas quedan apagadas»}. \textbf{Tres, no diez}: una regla que se cumple vale más que cinco que se rompen el jueves.}
\milcpaso{4}{milcMagenta}{\textbf{Poner el fuego (8 min).} Elige \textbf{un rato del día} en que va a haber un solo fuego y va a ser compartido: la comida, la caminada al colegio, media hora de juego con alguien, el rato antes de dormir con un libro o una conversación. \textbf{Escríbelo con hora y con quién.} \emph{La tulpa no era solo apagar lo demás: era estar mirando el mismo fuego que otro.}}

\begin{drawbox}{milcMagenta}{Antes de cerrar tu mapa de atención}
\begin{checklista}
\item Tienes tus números de ayer: tiempo total, tres apps y desbloqueos.
\item Identificaste al menos \textbf{tres enganches} concretos en apps que usas, señalándolos.
\item Tus tres reglas tienen \textbf{hora y lugar}, y son posibles en tu casa real.
\item Al menos una regla es un \textbf{ajuste que ya hiciste} hoy (apagar autoplay, apagar notificaciones), no una intención.
\item Escribiste tu fuego: qué rato, con quién y a qué hora.
\item Sabes cómo vas a medir dentro de una semana si cambió algo (los mismos tres datos).
\end{checklista}
\end{drawbox}

\begin{softbox}{milcMagenta}{milcGris}{Plantilla de trabajo}
\textbf{Plantilla del mapa.} \emph{Ayer:} \_\_\_\_ horas · apps: \_\_\_\_, \_\_\_\_, \_\_\_\_ · desbloqueos: \_\_\_\_ · elegido \_\_\_\_ / simplemente pasó \_\_\_\_. \\[1mm]
\textbf{Enganches que reconocí:} app \_\_\_\_ → enganche \_\_\_\_ ; app \_\_\_\_ → enganche \_\_\_\_. \\[1mm]
\textbf{Mis tres reglas:} 1) \_\_\_\_ (hora y lugar) · 2) \_\_\_\_ · 3) \_\_\_\_. \\[1mm]
\textbf{Mi fuego:} \emph{«\_\_\_\_ (qué), a las \_\_\_\_, con \_\_\_\_»}.
\end{softbox}

\begin{infoband}{milcMagenta}


\textbf{En tu cuaderno (Actividad 3 · CREA).} Título: «Actividad 3. Mi mapa de atención». \textbf{Formato:} los datos de ayer + los enganches reconocidos + las tres reglas con hora y lugar + el fuego elegido. \textbf{Extensión:} una página. \textbf{Sabes que terminaste cuando} \textbf{ya hiciste} al menos un ajuste en el teléfono, ahí mismo, y no solo lo escribiste.
\end{infoband}

\puente{El producto son tres reglas concretas y un fuego elegido. \emph{«Voy a usar menos el celular» no es una regla: es un deseo con buena intención.}}

\milcestacion{milcMostaza}{Producto de la sesión}
\textbf{Mapa de atención: datos de ayer, enganches reconocidos, tres reglas con hora y un fuego elegido}.
\begin{softbox}{milcMostaza}{milcGris}{Criterios de calidad}
(1) Los datos de ayer son datos: tiempo, aplicaciones y desbloqueos, no impresiones. (2) Se nombran al menos tres enganches concretos, cada uno en una aplicación específica. (3) Las tres reglas tienen hora y lugar y son posibles en la casa real de quien las escribe. (4) Al menos un ajuste está hecho, no prometido. (5) El fuego elegido dice qué, cuándo y con quién, y se puede sostener un martes normal.
\end{softbox}

\puente{Antes de cerrar, mira si tus reglas sobreviven a un martes normal. Si solo funcionan un domingo tranquilo, no son reglas.}

\milcestacion{milcVino}{Evaluación liberadora · cinco dimensiones}

\begin{softbox}{milcVino}{milcGris}{Sentido de la evaluación}
Evaluar no es solo revisar una nota. Es reconocer qué aprendí, qué cambió en mí, qué emociones manejé, cómo me relaciono con la ciudad y la comunidad, y qué le dejaré a quien viene detrás.
\end{softbox}

\textbf{Mi reflexión liberadora en cinco dimensiones.} Completa cada frase iniciada:

\small
\begin{tabularx}{\textwidth}{>{\bfseries\RaggedRight\arraybackslash}p{3.4cm}Y>{\RaggedRight\arraybackslash}p{4.6cm}}
\rowcolor{milcVino}\color{white}Dimensión & \color{white}\bfseries Mi respuesta (frame starter) & \color{white}\bfseries Algo que descubrí\\
1. Personal & Lo que más cambió en mí fue \linead{6.5cm} & \linead{4.2cm}\\
2. Emocional & La emoción más fuerte fue \linead{2.5cm} y la regulé \linead{3cm} & \linead{4.2cm}\\
3. Ciudadana & En democracia esto importa porque \linead{6cm} & \linead{4.2cm}\\
4. Local (Cartago) & En mi barrio sería útil para \linead{6.5cm} & \linead{4.2cm}\\
5. Intergeneracional & A mi abuela/abuelo le diría \linead{6.5cm} & \linead{4.2cm}\\
\end{tabularx}
\normalsize

\begin{infoband}{milcVino}
\textbf{Para la próxima sustentación.} Vuelve a esta tabla en seis meses. Verás que las respuestas cambian — no porque lo aprendido cambie, sino porque tú cambias. La evaluación liberadora es una conversación contigo a través del tiempo.
\end{infoband}


\puente{Cerramos con tres voces: quién queda por fuera cuando todos miran su pantalla, qué está en tu poder, y qué le hace al mundo compartido lo que consumes.}

\milcestacion{milcGradeDark}{Triángulo de pensamiento · cierre filosófico}

\begin{softbox}{milcVino}{milcGris}{Dussel · lente del nosotros}
\textit{«Aquel que es capaz de escuchar silenciosamente al Otro como otro es el que puede constituir una comunidad y no una mera sociedad totalizada.»}

{\footnotesize\itshape\color{milcNegro!70}--- Enrique Dussel · Filosofía de la liberación (1977)}

Alrededor de la tulpa se aprende \textbf{mirando y escuchando}: la atención puesta en el otro es lo que sostiene la comunidad. Cuando cuatro personas comparten una sala y cada una mira su propia luz, hay \emph{sociedad} ---están en el mismo lugar--- pero puede no haber \emph{comunidad}: nadie está escuchando a nadie. \textbf{Dar atención es lo más concreto que se le puede dar a alguien}, porque es limitada y no vuelve. \emph{Por eso mirar a alguien a la cara mientras te habla no es cortesía: es el material del que están hechos los vínculos.}

\textbf{¿A quién de mi casa hace cuánto no le doy diez minutos de atención completa, sin pantalla en la mano?}
\end{softbox}
\linead{\linewidth}\\[1mm]
\linead{\linewidth}

\begin{softbox}{milcMostaza}{milcGris}{Estoicismo · Epicteto · Enquiridión, 1 (c. 125 d.C.; trad. desde E. Carter) · cuidado interior}
\textit{«Algunas cosas están en nuestro poder y otras no. Están en nuestro poder la opinión, el impulso, el deseo, el rechazo; en una palabra, todo aquello que es obra nuestra.»}

{\footnotesize\itshape\color{milcNegro!70}--- Epicteto · Enquiridión, 1 (c. 125 d.C.; trad. desde E. Carter)}

No está en tu poder que las aplicaciones estén diseñadas para retenerte: eso lo deciden empresas grandes y no te preguntaron. \textbf{Sí está en tu poder dónde queda el teléfono mientras estudias, cuáles notificaciones dejas prendidas y a qué hora se apaga.} Fíjate en la trampa: si crees que todo depende de tu fuerza de voluntad, cada resbalón se siente como un defecto personal. \emph{Y si crees que nada depende de ti, no cambias nada. La verdad está en el medio, y la parte del medio es exactamente donde hay que trabajar.}

\textbf{De todo lo que me molesta de mi relación con el celular, ¿qué parte puedo cambiar hoy con un ajuste?}
\end{softbox}
\linead{\linewidth}\\[1mm]
\linead{\linewidth}

\begin{softbox}{milcTurquesa}{milcGris}{Floridi · lente de la infoesfera}
\textit{«Las tecnologías digitales están re-ontologizando la naturaleza misma de la infoesfera.»}

{\footnotesize\itshape\color{milcNegro!70}--- Luciano Floridi · La cuarta revolución (2014; trad.)}

\emph{Re-ontologizar} quiere decir que estas tecnologías no solo se meten en el mundo: \textbf{rehacen cómo es}. La sala de tu casa a las ocho de la noche no es la misma sala de hace veinte años; el descanso, la amistad y hasta el aburrimiento cambiaron de forma. \textbf{No es nostalgia ---nadie quiere volver atrás---, es darse cuenta de que el cambio no fue neutral} y de que alguien lo diseñó con intereses. \emph{Elegir tu fuego es participar en ese diseño en vez de solo recibirlo.}

\textbf{Si alguien mirara mis últimas 24 horas, ¿diría que yo elegí mi día o que me lo eligieron?}
\end{softbox}
\linead{\linewidth}\\[1mm]
\linead{\linewidth}

\begin{infoband}{milcNegro}
\textbf{Nota para el docente.} Las tres son citas textuales y llevan su referencia debajo. Puedes remitir a la obra en clase.
\end{infoband}

\milcestacion{milcVino}{Mi autoevaluación · marca una casilla por fila}

{\small
\setlength{\extrarowheight}{2.2pt}
\begin{tabularx}{\textwidth}{>{\RaggedRight\arraybackslash\bfseries}p{3.0cm}YYY}
\rowcolor{milcVino}\color{white}\bfseries Criterio & \color{white}\bfseries Lo logré & \color{white}\bfseries En proceso & \color{white}\bfseries Necesito apoyo\\[.6mm]
Comprensión del tema & \milcopcion{Explico las ideas con mis palabras.} & \milcopcion{Reconozco las ideas, falta precisión.} & \milcopcion{Necesito repasar lo básico.}\\[.8mm]
\rowcolor{milcGris}Lectura del enganche & \milcopcion{Reconozco en apps concretas los mecanismos que buscan quedarse con mi atención.} & \milcopcion{Reconozco que pierdo tiempo, pero no identifico qué lo produce.} & \milcopcion{Todavía creo que es solo falta de voluntad mía.}\\[.8mm]
Aplicación y producto & \milcopcion{Mi producto cumple los criterios.} & \milcopcion{Está armado, con ajustes pendientes.} & \milcopcion{Aún está en construcción.}\\[.8mm]
\rowcolor{milcGris}Reflexión 5 dimensiones & \milcopcion{Respondí cada una con honestidad.} & \milcopcion{Respondí algunas con detalle.} & \milcopcion{Mis respuestas son muy generales.}\\[.8mm]
Triángulo de pensamiento & \milcopcion{Conecté las 3 voces con mi trabajo.} & \milcopcion{Conecté al menos una.} & \milcopcion{No las relaciono con mi proceso.}\\[.8mm]
\end{tabularx}}

\vspace{3mm}
\textbf{Hoy entendí que:}\\[1mm]
\linead{\linewidth}\\[1mm]
\linead{\linewidth}

\textbf{Mi compromiso para la próxima sesión es:}\\[1mm]
\linead{\linewidth}\\[1mm]
\linead{\linewidth}

\begin{softbox}{milcMostaza}{milcGris}{Cita de despedida · Marco Aurelio}
\textit{``El obstáculo en el camino se vuelve el camino. Lo que estorba al camino se vuelve camino.''}\\[1mm]
Cada error de hoy, bien observado, se convierte en pieza del camino que pisarás mañana.
\end{softbox}

\small
\begin{tabularx}{\textwidth}{>{\bfseries}p{3.6cm}Y}
\rowcolor{milcGradeDark!12}Nombre completo: & \linead{10cm}\\
Fecha: & \linead{4cm} \quad Curso/grupo: \linead{4cm}\\
Mi firma: & \linead{9cm}\\
\end{tabularx}
\normalsize

\begin{infoband}{milcGradeDark}
\textbf{Para la próxima sesión.} Dos cosas. \textbf{Primero, cumple tus tres reglas una semana y vuelve a medir} los mismos tres datos ---tiempo, apps, desbloqueos---. \emph{Anota también las veces que la regla se rompió y qué la rompió: eso enseña más que el promedio.} \textbf{Segundo, prende tu fuego y cuéntalo:} sostén el rato que elegiste ---la comida sin pantallas, la caminada, la conversación--- \textbf{al menos tres veces}, y anota qué pasó: si costó, si alguien más se sumó, si alguien se molestó, de qué hablaron. \textbf{Trae:} tus números nuevos y una frase sobre cómo fue tu fuego. \emph{Ojo con algo: es probable que en tu casa haya adultos más enganchados que tú. Esto no es para reclamarles; si acaso, es para invitarlos.}
\end{infoband}

\end{document}
